% GENERATED FILE. Edit canonical lessons, then run npm run generate:books.
% canonical-source-hash: fnv1a64:f49455be6eb69567
% canonical-lessons: UR-C417-akaunt, UR-C417-raqam, UR-C417-naqd, UR-C417-chek, UR-C417-nikalvana, UR-R417-review-410-417

\chapter{Accounts and Cash}
\label{ch:ur-accounts-and-cash}
\hlchaptermodality{417}

\begin{chapteropening}
\textbf{By the end of this chapter:} \emph{I can talk about my bank account, cash and cheques.}

Complete the last lesson of chapter 417: I can talk about my bank account, cash and cheques.
\end{chapteropening}

\section[\texorpdfstring{\ur{اکاؤنٹ} (akāunṭ)}{اکاؤنٹ (akāunṭ)}]{\ur{اکاؤنٹ} (akāunṭ) — a bank account}

\label{lesson:UR-C417-akaunt}



\begin{quote}
Before the new one: say the Urdu for a stink, then the Urdu for to spread (something) around.
\end{quote}

\subsection*{You'll want to know: \ur{اکاؤنٹ}}
\textbf{\ur{اکاؤنٹ}} — \emph{akāunṭ} — \textquotedblleft{}a bank account\textquotedblright{}.

\begin{quote}
\ur{میرا} \ur{اکاؤنٹ} \ur{اس} \ur{بینک} \ur{میں} \ur{ہے۔}
\end{quote}

\begin{quote}
\emph{merā akāunṭ is baink meṅ hai.}
\end{quote}

\begin{quote}
My account is with this bank.
\end{quote}

Masculine.

\begin{cousinweb}
\textbf{\ur{اکاؤنٹ}} (\emph{akāunṭ}) is borrowed from English \emph{account}.
\end{cousinweb}

\begin{grammarlens}[title={What you've built}]
One more word for getting around.
\end{grammarlens}

\subsection*{Your turn}
\begin{itemize}
\raggedright
  \item \emph{Say it:} \emph{akāunṭ}
  \item \emph{Say it:} \emph{akāunṭ}, once more
  \item \emph{Say it:} \emph{merā akāunṭ is baink meṅ hai.} — My account is with this bank.
  \item \emph{Recall:} say the Urdu for a stink, then the Urdu for to spread (something) around, then say \emph{akāunṭ} again
\end{itemize}

\begin{tcolorbox}[breakable,colback=teal!4,colframe=teal!35!black,title={Before you move on}]
What does \ur{اکاؤنٹ} mean? (\textquotedblleft{}A bank account\textquotedblright{}.) Say it once more.
\end{tcolorbox}
\section[\texorpdfstring{\ur{رقم} (raqam)}{رقم (raqam)}]{\ur{رقم} (raqam) — a sum of money; an amount}

\label{lesson:UR-C417-raqam}



\begin{quote}
Before the new one: say the Urdu for to spread (something) around, then the Urdu for a bank account.
\end{quote}

\subsection*{You'll want to know: \ur{رقم}}
\textbf{\ur{رقم}} — \emph{raqam} — \textquotedblleft{}a sum of money; an amount\textquotedblright{}.

\begin{quote}
\ur{اکاؤنٹ} \ur{میں} \ur{رقم} \ur{کم} \ur{ہے۔}
\end{quote}

\begin{quote}
\emph{akāunṭ meṅ raqam kam hai.}
\end{quote}

\begin{quote}
There isn't much money in the account.
\end{quote}

Feminine.

\begin{cousinweb}
\textbf{\ur{رقم}} (\emph{raqam}) is usually traced, through Persian, to Arabic \emph{raqam}, \textquotedblleft{}a figure, a written mark\textquotedblright{}; in Urdu it is a sum of money.
\end{cousinweb}

\begin{grammarlens}[title={What you've built}]
One more word for getting around.
\end{grammarlens}

\subsection*{Your turn}
\begin{itemize}
\raggedright
  \item \emph{Say it:} \emph{raqam}
  \item \emph{Say it:} \emph{raqam}, once more
  \item \emph{Say it:} \emph{akāunṭ meṅ raqam kam hai.} — There isn't much money in the account.
  \item \emph{Recall:} say the Urdu for to spread (something) around, then the Urdu for a bank account, then say \emph{raqam} again
\end{itemize}

\begin{tcolorbox}[breakable,colback=teal!4,colframe=teal!35!black,title={Before you move on}]
What does \ur{رقم} mean? (\textquotedblleft{}A sum of money; an amount\textquotedblright{}.) Say it once more.
\end{tcolorbox}
\section[\texorpdfstring{\ur{نقد} (naqd)}{نقد (naqd)}]{\ur{نقد} (naqd) — cash}

\label{lesson:UR-C417-naqd}



\begin{quote}
Before the new one: say the Urdu for a bank account, then the Urdu for a sum of money.
\end{quote}

\subsection*{You'll want to know: \ur{نقد}}
\textbf{\ur{نقد}} — \emph{naqd} — \textquotedblleft{}cash\textquotedblright{}.

\begin{quote}
\ur{دکاندار} \ur{نقد} \ur{پیسے} \ur{مانگتا} \ur{ہے۔}
\end{quote}

\begin{quote}
\emph{dukāndār naqd paise māṅgtā hai.}
\end{quote}

\begin{quote}
The shopkeeper wants cash.
\end{quote}

Masculine.

\begin{cousinweb}
\textbf{\ur{نقد}} (\emph{naqd}) is usually traced, through Persian, to Arabic \emph{naqd}, \textquotedblleft{}ready money\textquotedblright{}.
\end{cousinweb}

\begin{grammarlens}[title={What you've built}]
One more word for getting around.
\end{grammarlens}

\subsection*{Your turn}
\begin{itemize}
\raggedright
  \item \emph{Say it:} \emph{naqd}
  \item \emph{Say it:} \emph{naqd}, once more
  \item \emph{Say it:} \emph{dukāndār naqd paise māṅgtā hai.} — The shopkeeper wants cash.
  \item \emph{Recall:} say the Urdu for a bank account, then the Urdu for a sum of money, then say \emph{naqd} again
\end{itemize}

\begin{tcolorbox}[breakable,colback=teal!4,colframe=teal!35!black,title={Before you move on}]
What does \ur{نقد} mean? (\textquotedblleft{}Cash\textquotedblright{}.) Say it once more.
\end{tcolorbox}
\section[\texorpdfstring{\ur{چیک} (chek)}{چیک (chek)}]{\ur{چیک} (chek) — a cheque}

\label{lesson:UR-C417-chek}



\begin{quote}
Before the new one: say the Urdu for a sum of money, then the Urdu for cash.
\end{quote}

\subsection*{You'll want to know: \ur{چیک}}
\textbf{\ur{چیک}} — \emph{chek} — \textquotedblleft{}a cheque\textquotedblright{}.

\begin{quote}
\ur{میں} \ur{نے} \ur{کرایے} \ur{کے} \ur{لیے} \ur{چیک} \ur{دیا۔}
\end{quote}

\begin{quote}
\emph{maiṅ ne kirāye ke liye chek diyā.}
\end{quote}

\begin{quote}
I paid the rent by cheque.
\end{quote}

Masculine.

\begin{cousinweb}
\textbf{\ur{چیک}} (\emph{chek}) is borrowed from English \emph{cheque}.
\end{cousinweb}

\begin{grammarlens}[title={What you've built}]
One more word for getting around.
\end{grammarlens}

\subsection*{Your turn}
\begin{itemize}
\raggedright
  \item \emph{Say it:} \emph{chek}
  \item \emph{Say it:} \emph{chek}, once more
  \item \emph{Say it:} \emph{maiṅ ne kirāye ke liye chek diyā.} — I paid the rent by cheque.
  \item \emph{Recall:} say the Urdu for a sum of money, then the Urdu for cash, then say \emph{chek} again
\end{itemize}

\begin{tcolorbox}[breakable,colback=teal!4,colframe=teal!35!black,title={Before you move on}]
What does \ur{چیک} mean? (\textquotedblleft{}A cheque\textquotedblright{}.) Say it once more.
\end{tcolorbox}
\section[\texorpdfstring{\ur{نکلوانا} (nikalvānā)}{نکلوانا (nikalvānā)}]{\ur{نکلوانا} (nikalvānā) — to withdraw (money); to have (something) taken out}

\label{lesson:UR-C417-nikalvana}



\begin{quote}
Before the new one: say the Urdu for cash, then the Urdu for a cheque.
\end{quote}

\subsection*{You'll want to know: \ur{نکلوانا}}
\textbf{\ur{نکلوانا}} — \emph{nikalvānā} — \textquotedblleft{}to withdraw (money); to have (something) taken out\textquotedblright{}.

\begin{quote}
\ur{سارا} \ur{نے} \ur{بینک} \ur{سے} \ur{پیسے} \ur{نکلوائے۔}
\end{quote}

\begin{quote}
\emph{sārā ne baink se paise nikalvāe.}
\end{quote}

\begin{quote}
Sara took some money out of the bank.
\end{quote}

Past, with \ur{نے} (he or she …): \ur{اس} \ur{نے} \ur{رقم} \ur{نکلوائی}, \emph{us ne raqam nikalvāī}.

\begin{cousinweb}
\textbf{\ur{نکلوانا}} (\emph{nikalvānā}) is built on \textbf{\ur{نکلنا}}, \textquotedblleft{}to come out\textquotedblright{}, with the ending \emph{-vā-} for getting someone else to do it: you have the bank take your money out. \textbf{\ur{نکلنا}} comes from \emph{niṣkalati}, a form CDIAL reconstructs (CDIAL 7478).
\end{cousinweb}

\begin{grammarlens}[title={What you've built}]
One more word for getting around.
\end{grammarlens}

\subsection*{Your turn}
\begin{itemize}
\raggedright
  \item \emph{Say it:} \emph{nikalvānā}
  \item \emph{Say it:} \emph{nikalvānā}, once more
  \item \emph{Say it:} \emph{sārā ne baink se paise nikalvāe.} — Sara took some money out of the bank.
  \item \emph{Recall:} say the Urdu for cash, then the Urdu for a cheque, then say \emph{nikalvānā} again
\end{itemize}

\begin{tcolorbox}[breakable,colback=teal!4,colframe=teal!35!black,title={Before you move on}]
What does \ur{نکلوانا} mean? (\textquotedblleft{}To withdraw (money); to have (something) taken out\textquotedblright{}.) Say it once more.
\end{tcolorbox}
\section[(dialogue)]{Looking back: from Landmarks to Accounts and Cash}

\label{lesson:UR-R417-review-410-417}



\begin{quote}
Say the Urdu for a building and a spot.
\end{quote}

\subsection*{Your turn}
Say these aloud:

\begin{itemize}
\raggedright
  \item a building, a spot, a signboard, a pole, to show
  \item a slope, to roll (down), a veranda, a basement, beside
  \item a small town, a population, an inhabitant, the countryside, next door
  \item a large house, a hut, a compound, a thatched roof, to found (a town)
  \item a guard, stray, to bark, to growl, to chase off
  \item a roadside eatery, a handcart, a betel leaf, a flavour, to serve (food)
  \item rubbish, a drain, a sewer, a stink, to spread (something) around
  \item a bank account, a sum of money, cash, a cheque, to withdraw (money)
\end{itemize}

\begin{tcolorbox}[breakable,colback=teal!4,colframe=teal!35!black,title={Before you move on}]
A building? (\textbf{\ur{عمارت}}, \emph{imārat}.) A spot? (\textbf{\ur{مقام}}, \emph{maqām}.) A signboard? (\textbf{\ur{بورڈ}}, \emph{borḍ}.) A pole? (\textbf{\ur{کھمبا}}, \emph{khambā}.) To show? (\textbf{\ur{دکھانا}}, \emph{dikhānā}.) A slope? (\textbf{\ur{ڈھلان}}, \emph{ḍhalān}.) To roll (down)? (\textbf{\ur{لڑھکنا}}, \emph{luṛhaknā}.) A veranda? (\textbf{\ur{برآمدہ}}, \emph{barāmadā}.) A basement? (\textbf{\ur{تہہ} \ur{خانہ}}, \emph{tah khānā}.) Beside? (\textbf{\ur{بغل}}, \emph{baghal}.) A small town? (\textbf{\ur{قصبہ}}, \emph{qasbā}.) A population? (\textbf{\ur{آبادی}}, \emph{ābādī}.) An inhabitant? (\textbf{\ur{باشندہ}}, \emph{bāshindā}.) The countryside? (\textbf{\ur{دیہات}}, \emph{dehāt}.) Next door? (\textbf{\ur{پڑوس}}, \emph{paṛos}.) A large house? (\textbf{\ur{کوٹھی}}, \emph{koṭhī}.) A hut? (\textbf{\ur{جھونپڑی}}, \emph{jhoṅpṛī}.) A compound? (\textbf{\ur{احاطہ}}, \emph{ahātā}.) A thatched roof? (\textbf{\ur{چھپر}}, \emph{chhappar}.) To found (a town)? (\textbf{\ur{بسانا}}, \emph{basānā}.) A guard? (\textbf{\ur{گارڈ}}, \emph{gārḍ}.) Stray? (\textbf{\ur{آوارہ}}, \emph{āvārā}.) To bark? (\textbf{\ur{بھونکنا}}, \emph{bhauṅknā}.) To growl? (\textbf{\ur{غرانا}}, \emph{ghurrānā}.) To chase off? (\textbf{\ur{بھگانا}}, \emph{bhagānā}.) A roadside eatery? (\textbf{\ur{ڈھابا}}, \emph{ḍhābā}.) A handcart? (\textbf{\ur{ٹھیلا}}, \emph{ṭhelā}.) A betel leaf? (\textbf{\ur{پان}}, \emph{pān}.) A flavour? (\textbf{\ur{ذائقہ}}, \emph{zāiqā}.) To serve (food)? (\textbf{\ur{پروسنا}}, \emph{parosnā}.) Rubbish? (\textbf{\ur{کوڑا}}, \emph{kūṛā}.) A drain? (\textbf{\ur{نالی}}, \emph{nālī}.) A sewer? (\textbf{\ur{گٹر}}, \emph{gaṭar}.) A stink? (\textbf{\ur{بدبو}}, \emph{badbū}.) To spread (something) around? (\textbf{\ur{پھیلانا}}, \emph{phailānā}.) A bank account? (\textbf{\ur{اکاؤنٹ}}, \emph{akāunṭ}.) A sum of money? (\textbf{\ur{رقم}}, \emph{raqam}.) Cash? (\textbf{\ur{نقد}}, \emph{naqd}.) A cheque? (\textbf{\ur{چیک}}, \emph{chek}.) To withdraw (money)? (\textbf{\ur{نکلوانا}}, \emph{nikalvānā}.)
\end{tcolorbox}
